\documentclass[12pt]{article}

\usepackage[margin=1in]{geometry}
\usepackage{newtxtext,newtxmath}
\usepackage{setspace}
\usepackage{graphicx}
\usepackage{booktabs}
\usepackage{tabularx}
\usepackage{array}
\usepackage{float}
\usepackage{xcolor}
\usepackage{hyperref}
\usepackage{enumitem}
\usepackage{caption}
\usepackage{fancyhdr}
\usepackage{microtype}
\usepackage{url}

% ------------------------------------------------------------
% Document style (reused from Midterm Project Report)
% ------------------------------------------------------------
\definecolor{linkblue}{HTML}{1A4F8B}

\hypersetup{
    colorlinks=true,
    linkcolor=linkblue,
    urlcolor=linkblue,
    citecolor=linkblue,
    pdftitle={Senior Seminar I Progress Report 3},
    pdfauthor={Musa Abdulrafiu}
}

\setstretch{1.15}
\setlength{\parindent}{0.5in}
\setlength{\parskip}{0.35em}
\setlist[itemize]{topsep=0.25em,itemsep=0.15em}
\captionsetup{font=small,labelfont=bf}

\pagestyle{fancy}
\fancyhf{}
\lhead{Senior Seminar I}
\rhead{Progress Report 3}
\cfoot{\thepage}
\renewcommand{\headrulewidth}{0.4pt}

% ------------------------------------------------------------
% Student/project information
% ------------------------------------------------------------
\newcommand{\SubmissionDate}{[FILL IN TODAY'S SUBMISSION DATE]}

\newcommand{\RepoURL}{https://github.com/MMAbdul153/agent-benchmark-python}

\newcommand{\StudentName}{Musa Abdulrafiu}
\newcommand{\InstructorName}{Dr.\ Ning Zhang}
\newcommand{\CourseName}{Senior Seminar I}

\newcommand{\AgentName}{Claude Code in VS Code}
\newcommand{\LLMName}{ChatGPT browser chat}

% ------------------------------------------------------------
% Screenshot placeholder / auto-image command.
%
% \ScreenshotPlaceholder{description}{filename-without-extension} checks
% whether images/<filename>.png already exists:
%   - If it does, it's included automatically as the figure -- no further
%     edits to this .tex file needed.
%   - If it doesn't (yet), a labeled dashed box is shown instead, naming
%     exactly what to capture and what to save it as.
%
% So the workflow is simply: save a screenshot to images/<filename>.png
% using the exact suggested name, recompile, and it appears in place of the
% box automatically. Nothing else in this file needs to change.
% ------------------------------------------------------------
\newcommand{\ScreenshotPlaceholderBox}[2]{%
  \fbox{\begin{minipage}[c][0.22\textheight][c]{0.92\textwidth}
    \centering
    \vspace{0.5em}
    {\large\bfseries [ SCREENSHOT NEEDED ]}\\[0.6em]
    #1\\[0.6em]
    {\small\ttfamily Suggested filename: images/#2.png}
    \vspace{0.5em}
  \end{minipage}}%
}
\newcommand{\ScreenshotPlaceholder}[2]{%
  \IfFileExists{images/#2.png}{%
    \includegraphics[width=0.95\textwidth]{images/#2.png}%
  }{%
    \ScreenshotPlaceholderBox{#1}{#2}%
  }%
}

\begin{document}

% ------------------------------------------------------------
% Title page (its own page, matching the Midterm Report's style)
% ------------------------------------------------------------
\begin{titlepage}
    \centering
    \vspace*{1.15in}

    {\Large \textbf{\CourseName}\par}
    \vspace{0.35in}

    {\LARGE \textbf{Progress Report 3}\par}
    \vspace{0.25in}

    {\Large \textbf{Benchmarking Autonomous AI Coding Agents Against Conventional LLM Code Generation}\par}
    \vspace{0.2in}

    {\large Post-Midterm Update: Weeks 8--9\par}

    \vfill

    \begin{tabular}{rl}
        \textbf{Student:} & \StudentName \\
        \textbf{Instructor:} & \InstructorName \\
        \textbf{Course:} & \CourseName \\
        \textbf{Submission Date:} & \SubmissionDate \\
        \textbf{GitHub Repository:} &
        \href{\RepoURL}{\texttt{github.com/MMAbdul153/agent-benchmark-python}} \\
    \end{tabular}

    \vfill
\end{titlepage}

% ------------------------------------------------------------
\section{Purpose of This Report}

The Midterm Project Report (Weeks 1--7) established the benchmark pipeline and
reported one completed experimental comparison: Task T01 (discount validation),
evaluated under both an agentic workflow (\AgentName) and a non-agentic workflow
(\LLMName), with both conditions passing 12 of 12 tests and 0 regressions.

This report documents work since the midterm on Task T02 (empty-list average
validation) and Task T03 (add-function test generation). Both tasks are now
fully closed out in both conditions (agentic and non-agentic), verified by
actually executing the test suite, using the same systems as T01: \AgentName{}
for the agentic condition and \LLMName{} for the non-agentic condition. In
keeping with the project's evidentiary standard, this report distinguishes
clearly between verified results and work that remains pending.

% ------------------------------------------------------------
\section{T02 Task Definition (Unchanged from Midterm Plan)}

Task T02 requires correcting the \texttt{average} function in
\texttt{src/calculator.py} so that:

\begin{itemize}
    \item A nonempty list returns its arithmetic average.
    \item A single-element list returns that element.
    \item An empty list raises \texttt{ValueError} (the original implementation
    incorrectly returned \texttt{0.0}).
    \item \texttt{apply\_discount} and \texttt{add} are left unmodified.
\end{itemize}

The expected baseline result, documented when the T02 task template was created
in Weeks 5--6, is exactly one failing test, \texttt{test\_empty\_list\_raises},
out of 12 total tests.

% ------------------------------------------------------------
\section{Work Completed and Verified: T02}

Following the same experimental-control pattern used for T01, isolated run
folders were created under \texttt{experimental\_runs/}:
\texttt{T02-agentic/} and \texttt{T02-nonagentic/}, each seeded from the
unmodified \texttt{task\_templates/T02\_average} starting state so that neither
condition's attempt can contaminate the other.

\textbf{Baseline (verified).} The unmodified T02 code was executed against the
12-test suite:

\begin{verbatim}
cd task_templates/T02_average
python -m unittest discover -s tests -p "test_calculator.py" -v
\end{verbatim}

Result: 12 tests run, exactly 1 failure (\texttt{test\_empty\_list\_raises}),
0 errors, matching the TASK.md specification. Full output saved to
\texttt{evidence/t02-baseline-output.txt}.

\begin{figure}[H]
\centering
\ScreenshotPlaceholder{T02 baseline test execution, showing 12 tests run and
1 failure (\texttt{test\_empty\_list\_raises}).}{t02-baseline}
\caption{T02 baseline test execution before any fix is applied.}
\label{fig:t02-baseline}
\end{figure}

\textbf{Agentic fix (verified).} The T02 agentic attempt was run with
\AgentName{}, scoped to only the \texttt{experimental\_runs/T02-agentic/}
folder. The agent changed
\texttt{average} to raise \texttt{ValueError("numbers must not be empty")} on
an empty list, leaving \texttt{apply\_discount}, \texttt{add}, and the test
file unchanged. It self-measured elapsed time using real start and end shell
timestamps (\texttt{2026-10-27T12:55:14} to \texttt{12:55:24}), yielding
\textbf{9.58 seconds} across \textbf{1}
conversational turn (3 tool calls). The final run produced 12/12 tests
passing, 0 regressions. The agent did not independently re-run the pre-fix
baseline; the starting file was byte-identical to the separately verified
\texttt{task\_templates/T02\_average} baseline (12 tests, 1 failure --
\texttt{test\_empty\_list\_raises}). Full output saved to
\texttt{evidence/t02-agentic-vscode-output.txt}. Note the 9.58-second figure
measures only agent-tool time between its own timestamp commands, not the
full human-observed wall-clock session used for T01's 45.10-second figure, so
the two are not strictly comparable either.

\begin{figure}[H]
\centering
\ScreenshotPlaceholder{Claude Code in VS Code, scoped to
experimental\_runs/T02-agentic/, showing the edited average() function
and/or the agent's reported timestamps.}{t02-agentic-session}
\caption{T02 agentic condition: Claude Code in VS Code editing
\texttt{average()}.}
\label{fig:t02-agentic-session}
\end{figure}

\begin{figure}[H]
\centering
\ScreenshotPlaceholder{Final T02 agentic test run, showing 12/12 tests
passing after the fix.}{t02-agentic-final}
\caption{T02 final test evidence, official \AgentName{} result.}
\label{fig:t02-agentic-final}
\end{figure}

\textbf{Non-agentic attempt (official, verified).} The static prompt packet
(\texttt{evidence/t02-nonagentic-prompt.txt}), following the same format used
for the T01 non-agentic condition with the real baseline output inserted, was
submitted to \LLMName. Its returned code was applied verbatim to
\texttt{experimental\_runs/T02-nonagentic/src/calculator.py} with no manual
correction: \texttt{average} now raises \texttt{ValueError("numbers cannot be
empty")} on an empty list, with \texttt{apply\_discount} and \texttt{add}
unchanged. Elapsed time was measured with an external stopwatch (prompt
submission to response finished rendering), the same method used for T01's
non-agentic condition: \textbf{9.03 seconds}, 1 turn. Re-running the test
suite confirmed 12/12 tests passing, 0 regressions. Full output saved to
\texttt{evidence/t02-nonagentic-output.txt}. Because this timing method
matches T01's non-agentic measurement, the 9.03s figure is directly
comparable to T01's 4.32s -- but neither is comparable to any agentic row,
which uses a different timing definition (see Section~1 and Section~3).

\begin{figure}[H]
\centering
\ScreenshotPlaceholder{ChatGPT browser chat: the submitted T02 non-agentic
prompt and its returned code.}{t02-nonagentic-chat}
\caption{T02 non-agentic condition: ChatGPT prompt and response.}
\label{fig:t02-nonagentic-chat}
\end{figure}

\begin{figure}[H]
\centering
\ScreenshotPlaceholder{Final T02 non-agentic test run, after applying
ChatGPT's returned code verbatim.}{t02-nonagentic-final}
\caption{T02 final test evidence, non-agentic (\LLMName) result.}
\label{fig:t02-nonagentic-final}
\end{figure}

% ------------------------------------------------------------
\section{Work Completed and Verified: T03}

The same pattern was applied to T03 (add-function test generation), with
isolated folders \texttt{experimental\_runs/T03-agentic/} and
\texttt{T03-nonagentic/} seeded from \texttt{task\_templates/T03\_add\_tests}.

\textbf{Baseline (verified).}

\begin{verbatim}
cd task_templates/T03_add_tests
python -m unittest discover -s tests -p "test_calculator.py" -v
\end{verbatim}

Result: 8 tests run, 0 failures, 0 errors — only \texttt{TestDiscount} and
\texttt{TestAverage} exist before this task, as specified. A passing baseline
does not indicate completion, since the required \texttt{TestAdd} tests are
the task itself. Full output saved to \texttt{evidence/t03-baseline-output.txt}.

\begin{figure}[H]
\centering
\ScreenshotPlaceholder{T03 baseline test execution, showing 8 tests run, 0
failures, and no TestAdd class present yet.}{t03-baseline}
\caption{T03 baseline test execution before any tests are added.}
\label{fig:t03-baseline}
\end{figure}

\textbf{Agentic test generation (verified).} The T03 agentic attempt was run
with \AgentName{}, scoped to only the
\texttt{experimental\_runs/T03-agentic/} folder. The agent added a
\texttt{TestAdd} class with four tests -- positive integers, floating-point
values, negative numbers (one negative plus one positive, and two negatives),
and large values in the billions -- leaving \texttt{add}, \texttt{apply\_discount},
\texttt{average}, and all existing tests unchanged. It self-reported reading
and editing files only inside its assigned run folder. Using real start/end
shell timestamps, elapsed time was \textbf{9.11 seconds} across \textbf{1}
conversational turn (4 tool calls). The final run produced 12/12 tests
passing, 0 regressions. As with T02, the agent did not independently re-run
the pre-change baseline; the starting file was byte-identical to the
separately verified \texttt{task\_templates/T03\_add\_tests} baseline (8
tests, 0 failures, no \texttt{TestAdd} present). Full output saved to
\texttt{evidence/t03-agentic-vscode-output.txt}.

\begin{figure}[H]
\centering
\ScreenshotPlaceholder{Claude Code in VS Code, scoped to
experimental\_runs/T03-agentic/, showing the added TestAdd class.}{t03-agentic-session}
\caption{T03 agentic condition: Claude Code in VS Code adding
\texttt{TestAdd}.}
\label{fig:t03-agentic-session}
\end{figure}

\begin{figure}[H]
\centering
\ScreenshotPlaceholder{Final T03 agentic test run, showing 12/12 tests
passing (8 existing + 4 new).}{t03-agentic-final}
\caption{T03 final test evidence, official \AgentName{} result.}
\label{fig:t03-agentic-final}
\end{figure}

\textbf{Non-agentic attempt (official, verified).} The corresponding prompt
packet (\texttt{evidence/t03-nonagentic-prompt.txt}, with the real baseline
output inserted) was submitted to \LLMName. Its returned test file was
applied with one correction: the pasted response read
\texttt{if name == "main":} instead of
\texttt{if \_\_name\_\_ == "\_\_main\_\_":} -- a markdown-paste artifact (chat
markdown's bold-text parsing strips double underscores), not a logic change,
so the literal dunders were restored rather than treated as part of the
model's output. No test logic or assertion was altered. The returned
\texttt{TestAdd} class (4 tests: positive integers, floats, negative numbers,
numbers at the \(10^{15}\) scale) was otherwise applied verbatim, with
\texttt{TestDiscount} and \texttt{TestAverage} unchanged. External stopwatch
time (same method as T01 and T02's non-agentic conditions): \textbf{5.28
seconds}, 1 turn. Re-running the test suite confirmed 12/12 tests passing, 0
regressions. Full output saved to \texttt{evidence/t03-nonagentic-output.txt}.

\begin{figure}[H]
\centering
\ScreenshotPlaceholder{ChatGPT browser chat: the submitted T03 non-agentic
prompt and its returned test file.}{t03-nonagentic-chat}
\caption{T03 non-agentic condition: ChatGPT prompt and response.}
\label{fig:t03-nonagentic-chat}
\end{figure}

\begin{figure}[H]
\centering
\ScreenshotPlaceholder{Final T03 non-agentic test run, after applying
ChatGPT's returned test file verbatim.}{t03-nonagentic-final}
\caption{T03 final test evidence, non-agentic (\LLMName) result.}
\label{fig:t03-nonagentic-final}
\end{figure}

% ------------------------------------------------------------
\section{Verification Status Summary}

\begin{table}[H]
\centering
\caption{T02 and T03 verification status, this reporting period}
\label{tab:t02-t03-status}
\small
\begin{tabularx}{\textwidth}{l l c c c X}
\toprule
\textbf{Task} & \textbf{Condition} & \textbf{Tests} & \textbf{Regressions} & \textbf{Verified} & \textbf{Notes} \\
\midrule
T02 & Agentic (\AgentName) & 12/12 & 0 & \textbf{Yes} & 9.58s, 1 turn; see Section~3 \\
T02 & Non-agentic & 12/12 & 0 & \textbf{Yes} & 9.03s (stopwatch, same method as T01) \\
T03 & Agentic (\AgentName) & 12/12 & 0 & \textbf{Yes} & 9.11s, 1 turn; see Section~4 \\
T03 & Non-agentic & 12/12 & 0 & \textbf{Yes} & 5.28s (stopwatch, same method as T01/T02) \\
\bottomrule
\end{tabularx}
\end{table}

All six cells above (T01--T03, agentic and non-agentic) use the same two
systems throughout: \AgentName{} for every agentic attempt and \LLMName{} for
every non-agentic attempt, so the agentic condition is directly comparable
across tasks.

\begin{figure}[H]
\centering
\ScreenshotPlaceholder{results.csv showing the T01 rows alongside the new,
fully verified T02/T03 rows.}{results-csv-update}
\caption{Updated structured result log after this reporting period.}
\label{fig:results-csv-update}
\end{figure}

% ------------------------------------------------------------
\section{Remaining Steps}

All four T02/T03 experimental cells (agentic and non-agentic, both tasks) are
now complete and verified. What remains:

\begin{enumerate}
    \item Capture terminal screenshots for the T02 and T03 baseline, official
    agentic, and non-agentic results, to be added as figures in the next
    report revision (placeholders already inserted below).
    \item Decide whether to run the SHA-256 isolation check
    (\small\texttt{evidence/t02-t03-isolation-check.ps1}\normalsize)
    retroactively or treat the scoped-VS-Code-workspace control as
    sufficient; note the choice explicitly as a limitation either way.
    \item Begin repeated trials for T01--T03 using the new automation
    scripts (\texttt{scripts/bench\_reset.py}, \texttt{scripts/bench\_log.py}),
    since every result so far is a single trial per condition per task.
\end{enumerate}

% ------------------------------------------------------------
\section{Updated Plan for the Remainder of the Semester}

\begin{enumerate}
    \item Complete the non-agentic conditions for T02 and T03 (steps above).
    \item Run repeated trials for both the agentic and non-agentic conditions
    across multiple tasks.
    \item Standardize timing definitions and measurement procedures across all
    conditions and systems, including the T02/T03 agentic system-difference
    noted in Section~1.
    \item Add a simple code-quality rubric addressing readability,
    maintainability, scope of change, and regressions.
    \item Analyze quantitative results and qualitative failure modes for the
    final report.
\end{enumerate}

% ------------------------------------------------------------
\section{Evidence Produced This Period}

\begin{itemize}
    \item \texttt{experimental\_runs/T02-agentic/} --- official \AgentName{}
    fix applied and verified (12/12 tests, 0 regressions, 9.58s, 1 turn).
    \item \texttt{experimental\_runs/T03-agentic/} --- official \AgentName{}
    TestAdd class added and verified (12/12 tests, 0 regressions, 9.11s,
    1 turn).
    \item \texttt{experimental\_runs/T02-nonagentic/} --- official \LLMName{}
    fix applied and verified (12/12 tests, 0 regressions, 9.03s, 1 turn).
    \item \texttt{experimental\_runs/T03-nonagentic/} --- official \LLMName{}
    TestAdd class applied and verified (12/12 tests, 0 regressions, 5.28s,
    1 turn).
    \item \texttt{evidence/t02-nonagentic-prompt.txt},
    \texttt{evidence/t03-nonagentic-prompt.txt} --- static prompt packets for
    the non-agentic conditions, with real baseline output inserted and a note
    directing external stopwatch timing instead of model self-report.
    \item \texttt{evidence/t02-nonagentic-output.txt},
    \texttt{evidence/t03-nonagentic-output.txt} --- official, verified T02 and
    T03 non-agentic results (\LLMName).
    \item \texttt{evidence/t02-baseline-output.txt},
    \texttt{evidence/t03-baseline-output.txt} --- verified baseline terminal
    output.
    \item \small\texttt{evidence/t02-agentic-vscode-output.txt},
    \texttt{evidence/t03-agentic-vscode-output.txt}\normalsize{} --- official,
    verified T02 and T03 agentic results (\AgentName).
    \item \texttt{evidence/t02-t03-isolation-check.ps1} --- SHA-256 isolation
    check script, run before/after each agentic session to confirm no file
    outside the active run folder changed.
    \item \texttt{results.csv}, \texttt{summary.csv} --- updated with verified
    rows for all six T02/T03 cells (agentic and non-agentic, both tasks).
\end{itemize}

Terminal-output screenshots are not yet included in this report; they will be
added once captured, alongside the non-agentic results, in the next revision.

\end{document}
